% =============================================================================
%  RL-H2-02 · Tabel positioning vs paper terdekat
% -----------------------------------------------------------------------------
%  Role    : RL — Research Lead
%  Sprint  : H2 (setelah RL-H1-02) · PARALEL
%  Ticket  : (isi dari papan kerja)          Status: STARTER ONLY
%  Target  : docs/positioning.tex  ->  docs/positioning.pdf
%  Input   : docs/lit/*.tex + abstrak + brief v2 §2/§6
%  Output  : dokumentasi pembeda
%  Handoff : RL-H4-01 (related work), laporan
%  DoD     : [ ] baris "studi ini" hanya berisi klaim AMAN
%            [ ] setiap sel pesaing bersumber (lit note/abstrak)
%            [ ] 3 kalimat posisi konsisten dengan batas klaim brief §6
% =============================================================================
\documentclass[11pt,a4paper]{article}
% =============================================================================
%  00_preamble.tex | shared preamble untuk semua deliverable RL
%  Job     : bukan job card; dipakai oleh RL-H1-01 .. RL-H7-01
%  Target  : docs/ (salin bersama file .tex)
%  Catatan : kompilasi dengan pdflatex. Untuk XeLaTeX/LuaLaTeX, hapus inputenc.
% =============================================================================

\usepackage[utf8]{inputenc}
\usepackage[T1]{fontenc}
\usepackage{lmodern}
\usepackage[a4paper,margin=2.4cm]{geometry}
\usepackage{amsmath,amssymb}
\usepackage{booktabs}
\usepackage{longtable}
\usepackage{array}
\usepackage{graphicx}
\usepackage{xcolor}
\usepackage{enumitem}
\usepackage{float}
\usepackage{caption}
\usepackage{url}
\usepackage{microtype}
\usepackage[colorlinks=true,linkcolor=blue!50!black,urlcolor=blue!50!black,citecolor=blue!50!black]{hyperref}

% --- helper wajib -----------------------------------------------------------
% Path/nama file aman (underscore, garis miring) di dalam teks.
\newcommand{\fpath}[1]{\texttt{\detokenize{#1}}}
% TODO yang mencolok. Jangan dihapus sebelum job selesai.
\newcommand{\todo}[1]{\textcolor{red!75!black}{\textbf{[TODO: #1]}}}
\newcommand{\jobid}[1]{\textsf{\textbf{#1}}}
\newcommand{\runid}[1]{\texttt{\detokenize{#1}}}
\newcommand{\claim}[1]{\textsf{#1}}
\newcommand{\good}{\textcolor{green!45!black}{$\checkmark$}}
\newcommand{\bad}{\textcolor{red!75!black}{$\times$}}
\newcommand{\maybe}{\textcolor{orange!80!black}{$\sim$}}

\newcolumntype{L}[1]{>{\raggedright\arraybackslash}p{#1}}
\newcolumntype{C}[1]{>{\centering\arraybackslash}p{#1}}

\setlength{\parindent}{0pt}
\setlength{\parskip}{4pt}
\captionsetup{font=small,labelfont=bf}
\setlist[itemize]{leftmargin=1.2em,itemsep=2pt,topsep=2pt}
\setlist[enumerate]{leftmargin=1.5em,itemsep=2pt,topsep=2pt}

% Judul job yang konsisten di semua dokumen.
\newcommand{\jobheader}[4]{%
  \begin{center}
  {\Large\bfseries #1}\\[2pt]
  {\large #2}\\[4pt]
  \begin{tabular}{@{}ll@{}}
  \toprule
  Role & #3 \\
  Sprint & #4 \\
  \bottomrule
  \end{tabular}
  \end{center}
}


\title{\jobid{RL-H2-02} | Positioning vs Paper Terdekat\\
\large Apa yang sudah ada, apa yang belum, dan klaim kita}
\author{RL \\ \small RENARO Lab}
\date{23 November 2026 \quad | \quad \textcolor{red!75!black}{STARTER}}

\begin{document}
\maketitle

\section{Properti graph}

\begin{table}[H]
\centering
\footnotesize
\setlength{\tabcolsep}{3pt}
\begin{tabular}{@{}L{2.2cm}L{4.0cm}C{1.5cm}C{1.5cm}L{2.5cm}L{2.7cm}@{}}
\toprule
Paper & Domain/data & Berarah & Dinamis & Jeda/edge & Kondisi graph \\
\midrule
RCSGNet & Kentang, 25 provinsi Tiongkok, mingguan 2012--2018
& \todo{} & \todo{} & \todo{} & laten (harga) \\
MarketDNA & Kentang multi-hub & \todo{} & \todo{} & \todo{} & laten diskret \\
MillGNN & multi-skala lead-lag & \todo{} & \todo{} & \todo{} & jendela input \\
SAMSGL & series-aligned & \todo{} & \todo{} & \todo{} & \todo{} \\
CRAFormer & longsor, driver hujan & \todo{} & \todo{} & \todo{} & agregat tetap \\
DRAN & spatio-temporal umum & \todo{} & \todo{} & \todo{} & gate statis/dinamis \\
SDGF & multi-skala & \todo{} & \todo{} & \todo{} & attention gate \\
STEVE & traffic & \todo{} & \todo{} & \todo{} & confounder laten \\
CGNN & harga pertanian & \todo{} & \todo{} & \todo{} & graph jarak statis \\
Min et al.\ 2025 & harga pertanian, Korea & \todo{} & \todo{} & \todo{} & \todo{} \\
Teh Sri Lanka & teh, jeda cuaca sumber & \todo{} & \todo{} & 1--3 minggu & tanpa graph \\
\midrule
\textbf{Studi ini} & Cabai merah, 34 provinsi Indonesia, harian kerja
& ya & ya & distribusi per edge & \textbf{iklim teramati asal} \\
\bottomrule
\end{tabular}
\caption{Properti graph. Sel \todo{} diisi dari full-text, bukan dari abstrak,
kecuali baris yang sudah ditandai di tabel sumber.}
\end{table}

\section{Properti evaluasi}

\begin{table}[H]
\centering
\footnotesize
\setlength{\tabcolsep}{4pt}
\begin{tabular}{@{}L{2.4cm}C{2.0cm}C{1.8cm}L{5.6cm}@{}}
\toprule
Paper & Iklim eksogen & Uji OOD regime & Uji graph prior keliru \\
\midrule
RCSGNet & \todo{} & \todo{} & \todo{} \\
MarketDNA & \todo{} & \todo{} & \todo{} \\
MillGNN & \todo{} & \todo{} & \todo{} \\
SAMSGL & \todo{} & \todo{} & \todo{} \\
CRAFormer & ya (hujan) & \todo{} & \todo{} \\
DRAN & \todo{} & \todo{} & \todo{} \\
SDGF & \todo{} & \todo{} & \todo{} \\
STEVE & tidak langsung & \todo{} & \todo{} \\
CGNN & ya (fitur node) & \todo{} & \todo{} \\
Min et al.\ 2025 & ya (node) & \todo{} & \todo{} \\
Teh Sri Lanka & ya (lag sumber) & \todo{} & \todo{} \\
\midrule
\textbf{Studi ini} & ya (asal) & ya & ya \\
\bottomrule
\end{tabular}
\caption{Properti evaluasi. Kolom OOD memisahkan uji regime iklim dari uji
weekday/holiday biasa.}
\end{table}

\section{Sumber setiap baris}
\begin{table}[H]
\centering
\small
\begin{tabular}{@{}L{3.2cm}L{6.0cm}L{5.2cm}@{}}
\toprule
Paper & Sumber & Halaman/bagian \\
\midrule
RCSGNet & \todo{\fpath{docs/lit/rcsgnet.tex}} & \todo{} \\
MarketDNA & \todo{} & \todo{} \\
MillGNN & \todo{} & \todo{} \\
SAMSGL & \todo{} & \todo{} \\
CRAFormer & \todo{} & \todo{} \\
DRAN & \todo{} & \todo{} \\
SDGF & \todo{} & \todo{} \\
STEVE & \todo{} & \todo{} \\
CGNN & \todo{} & \todo{} \\
Min et al.\ 2025 & \todo{} & \todo{} \\
Teh Sri Lanka & \todo{} & \todo{} \\
\bottomrule
\end{tabular}
\end{table}

\section{Tiga kalimat posisi (wajib konsisten dengan brief \S6)}
\begin{enumerate}
  \item Sepengetahuan kami, belum ada ST-GNN antarprovinsi untuk harga cabai
        Indonesia yang mengondisikan \emph{distribusi jeda per edge} pada iklim
        teramati di wilayah asal. \todo{verifikasi setelah full-text}
  \item Studi ini menyediakan benchmark evaluasi terstratifikasi regime dengan
        fitur \emph{availability-safe}, \emph{worst-regime error}, dan uji
        misspesifikasi graph prior. \todo{konfirmasi}
  \item Kami \textbf{tidak} mengklaim graph berarah/adaptif, gating, atau
        penggunaan cuaca pada GNN harga pertanian sebagai hal baru; edge adalah
        dependensi prediktif, bukan arus pasokan atau kausalitas.
\end{enumerate}

\section{Batas klaim yang harus dipatuhi}
\begin{table}[H]
\centering
\small
\begin{tabular}{@{}L{2.4cm}L{6.6cm}L{5.4cm}@{}}
\toprule
Kategori & Boleh diklaim & Tidak boleh \\
\midrule
AMAN & \todo{salin frasa brief \S6} & -- \\
BERISIKO & \todo{beri kualifikasi} & \todo{} \\
JANGAN & -- & graph adaptif/gating sebagai novelty; edge = arus/kausalitas;
``pertama memakai iklim pada GNN harga pertanian'' \\
\bottomrule
\end{tabular}
\end{table}

\end{document}
